\documentclass[12pt,a4paper]{exam}
\usepackage[utf8]{inputenc}
\usepackage{amsmath,amssymb,amsfonts}
\usepackage{geometry}
\usepackage{enumitem}
\usepackage{graphicx}
\usepackage{hyperref}
\usepackage{xcolor}
\geometry{a4paper, top=2cm, bottom=2cm, left=2.5cm, right=2.5cm}
\hypersetup{colorlinks=true, linkcolor=blue, urlcolor=blue}
\renewcommand{\thequestion}{\textbf{\arabic{question}}}
\renewcommand{\questionlabel}{\thequestion.}
\pointname{ marks}
\pointsdroppedatright
\bracketedpoints
\setlength{\parindent}{0pt}
\setlength{\emergencystretch}{3em}
\allowdisplaybreaks
\newsavebox{\schmathbox}
\newcommand{\fitmath}[1]{%
  \sbox{\schmathbox}{$\displaystyle #1$}%
  \ifdim\wd\schmathbox>\linewidth
    \resizebox{\linewidth}{!}{\usebox{\schmathbox}}%
  \else
    \usebox{\schmathbox}%
  \fi}

\begin{document}

\thispagestyle{empty}
\begin{center}
  \textbf{\LARGE Diag} \\[0.3cm]
  \textbf{Time Allowed: 3 Hours} \hfill \textbf{Maximum Marks: 80} \\[0.3cm]
  \rule{\textwidth}{0.5pt}
\end{center}

\vspace{0.3cm}

\noindent\textbf{General Instructions:}
\begin{enumerate}
  \item {\normalfont\normalcolor This question paper contains 40 questions. All questions are compulsory.}
  \item {\normalfont\normalcolor Questions are grouped into sections A through E.}
  \item {\normalfont\normalcolor Use of calculators is NOT permitted.}
\end{enumerate}
\bigskip
\noindent\textbf{\large Section A: Multiple Choice \& Assertion-Reason (1 mark each)}

\begin{questions}
\question[1] If $\cos \theta = 0$, then the value of $\theta$ (where $0^\circ \leq \theta \leq 90^\circ$) is:
\begin{enumerate}[label=(\Alph*)]
  \item $(A) 0^\circ$
  \item $(B) 45^\circ$
  \item $(C) 90^\circ$
  \item $(D) 30^\circ$
\end{enumerate}
\vspace{0.15cm}

\question[1] The value of $\sin^2 30^\circ + \cos^2 30^\circ$ is:
\begin{enumerate}[label=(\Alph*)]
  \item (A) 1
  \item (B) 0
  \item (C) 2
  \item (D) 1/2
\end{enumerate}
\vspace{0.15cm}

\question[1] If $\sec \theta = 2$, then $\theta$ is:
\begin{enumerate}[label=(\Alph*)]
  \item $(A) 30^\circ$
  \item $(B) 60^\circ$
  \item $(C) 45^\circ$
  \item $(D) 0^\circ$
\end{enumerate}
\vspace{0.15cm}

\question[1] The value of $\frac{\tan 30^\circ}{\cot 60^\circ}$ is:
\begin{enumerate}[label=(\Alph*)]
  \item (A) 1
  \item (B) 1/3
  \item (C) 3
  \item (D) 0
\end{enumerate}
\vspace{0.15cm}

\question[1] Which of the following is equal to $\cos \theta$?
\begin{enumerate}[label=(\Alph*)]
  \item $(A) \sin \theta$
  \item $(B) 1/\tan \theta$
  \item $(C) \sin(90^\circ + \theta)$
  \item $(D) \sin(90^\circ - \theta)$
\end{enumerate}
\vspace{0.15cm}

\question[1] If $3\tan \theta = 3$, then $\theta$ is:
\begin{enumerate}[label=(\Alph*)]
  \item $(A) 30^\circ$
  \item $(B) 45^\circ$
  \item $(C) 60^\circ$
  \item $(D) 90^\circ$
\end{enumerate}
\vspace{0.15cm}

\question[1] The value of $\sin 0^\circ + \cos 0^\circ$ is:
\begin{enumerate}[label=(\Alph*)]
  \item (A) 0
  \item (B) 2
  \item (C) 1
  \item (D) -1
\end{enumerate}
\vspace{0.15cm}

\question[1] If $\sin \theta = \cos \theta$, then $\theta$ is:
\begin{enumerate}[label=(\Alph*)]
  \item $(A) 30^\circ$
  \item $(B) 45^\circ$
  \item $(C) 60^\circ$
  \item $(D) 0^\circ$
\end{enumerate}
\vspace{0.15cm}

\question[1] If $\sin A = \frac{3}{5}$, then the value of $\cos A$ is:
\begin{enumerate}[label=(\Alph*)]
  \item (A) 2/5
  \item (B) 3/4
  \item (C) 4/5
  \item (D) 1/5
\end{enumerate}
\vspace{0.15cm}

\question[1] The value of $\tan 45^\circ$ is:
\begin{enumerate}[label=(\Alph*)]
  \item (A) 0
  \item (B) 1
  \item $(C) \sqrt{3}$
  \item $(D) 1/\sqrt{3}$
\end{enumerate}
\vspace{0.15cm}

\question[1] In a right triangle $ABC$ right-angled at $B$, if $AB=3$ cm and $BC=4$ cm, then $\sin A$ is:
\begin{enumerate}[label=(\Alph*)]
  \item (A) 4/5
  \item (B) 3/5
  \item (C) 3/4
  \item (D) 4/3
\end{enumerate}
\vspace{0.15cm}

\question[1] The value of $\cos^2 30^\circ + \sin^2 30^\circ$ is:
\begin{enumerate}[label=(\Alph*)]
  \item (A) 1/2
  \item (B) 0
  \item (C) 1
  \item (D) 2
\end{enumerate}
\vspace{0.15cm}

\question[1] If $\tan \theta = \frac{1}{\sqrt{3}}$, then the value of $\theta$ is:
\begin{enumerate}[label=(\Alph*)]
  \item $(A) 45^\circ$
  \item $(B) 30^\circ$
  \item $(C) 60^\circ$
  \item $(D) 90^\circ$
\end{enumerate}
\vspace{0.15cm}

\question[1] The value of $\frac{\sin 18^\circ}{\cos 72^\circ}$ is:
\begin{enumerate}[label=(\Alph*)]
  \item (A) 0
  \item (B) 1/2
  \item (C) 1
  \item $(D) \sqrt{3}$
\end{enumerate}
\vspace{0.15cm}

\question[1] If $3\cot A = 4$, then $\tan A$ is:
\begin{enumerate}[label=(\Alph*)]
  \item (A) 4/3
  \item (B) 3/5
  \item (C) 5/3
  \item (D) 3/4
\end{enumerate}
\vspace{0.15cm}

\question[1] The value of $9\sec^2 A - 9\tan^2 A$ is:
\begin{enumerate}[label=(\Alph*)]
  \item (A) 9
  \item (B) 1
  \item (C) 0
  \item (D) 8
\end{enumerate}
\vspace{0.15cm}

\question[1] If $\cos \theta = 0$, then the value of $\theta$ is:
\begin{enumerate}[label=(\Alph*)]
  \item $(A) 0^\circ$
  \item $(B) 90^\circ$
  \item $(C) 45^\circ$
  \item $(D) 60^\circ$
\end{enumerate}
\vspace{0.15cm}

\question[1] The value of $\sin 0^\circ + \cos 0^\circ$ is:
\begin{enumerate}[label=(\Alph*)]
  \item (A) 0
  \item (B) 1/2
  \item (C) 1
  \item (D) -1
\end{enumerate}
\vspace{0.15cm}

\question[1] Assertion (A): For an acute angle $\theta$, if $\sin \theta = \cos \theta$, then $\theta = 45^\circ$. 
Reason (R): The value of $\sin \theta$ is equal to $\cos \theta$ only at $\theta = 45^\circ$ in the interval $0^\circ \le \theta \le 90^\circ$.
\begin{enumerate}[label=(\Alph*)]
  \item (A) Both Assertion (A) and Reason (R) are true and Reason (R) is the correct explanation of Assertion (A)
  \item (B) Both Assertion (A) and Reason (R) are true but Reason (R) is NOT the correct explanation of Assertion (A)
  \item (C) Assertion (A) is true but Reason (R) is false
  \item (D) Assertion (A) is false but Reason (R) is true
\end{enumerate}
\vspace{0.15cm}

\question[1] Assertion (A): The value of $2 \sin^2 30^\circ - 3 \cos^2 45^\circ + 1$ is $0$. 
Reason (R): The value of $\sin 30^\circ = 1/2$ and $\cos 45^\circ = 1/\sqrt{2}$.
\begin{enumerate}[label=(\Alph*)]
  \item (A) Both Assertion (A) and Reason (R) are true and Reason (R) is the correct explanation of Assertion (A)
  \item (B) Both Assertion (A) and Reason (R) are true but Reason (R) is NOT the correct explanation of Assertion (A)
  \item (C) Assertion (A) is true but Reason (R) is false
  \item (D) Assertion (A) is false but Reason (R) is true
\end{enumerate}
\vspace{0.15cm}

\question[1] Assertion (A): The value of $\sin \theta$ can be $2$ for some angle $\theta$. 
Reason (R): The range of $\sin \theta$ is $[-1, 1]$.
\begin{enumerate}[label=(\Alph*)]
  \item (A) Both Assertion (A) and Reason (R) are true and Reason (R) is the correct explanation of Assertion (A)
  \item (B) Both Assertion (A) and Reason (R) are true but Reason (R) is NOT the correct explanation of Assertion (A)
  \item (C) Assertion (A) is true but Reason (R) is false
  \item (D) Assertion (A) is false but Reason (R) is true
\end{enumerate}
\vspace{0.15cm}

\question[1] Assertion (A): In a right-angled triangle, if $\tan A = 3/4$, then $\sin A = 3/5$. 
Reason (R): $\sin A = \frac{\text{Opposite}}{\text{Hypotenuse}}$ and $\tan A = \frac{\text{Opposite}}{\text{Adjacent}}$.
\begin{enumerate}[label=(\Alph*)]
  \item (A) Both Assertion (A) and Reason (R) are true and Reason (R) is the correct explanation of Assertion (A)
  \item (B) Both Assertion (A) and Reason (R) are true but Reason (R) is NOT the correct explanation of Assertion (A)
  \item (C) Assertion (A) is true but Reason (R) is false
  \item (D) Assertion (A) is false but Reason (R) is true
\end{enumerate}
\vspace{0.15cm}

\question[1] Assertion (A): The value of $\cos 0^\circ \cdot \cos 1^\circ \cdot \cos 2^\circ \dots \cos 90^\circ$ is $1$. 
Reason (R): $\cos 90^\circ = 0$.
\begin{enumerate}[label=(\Alph*)]
  \item (A) Both Assertion (A) and Reason (R) are true and Reason (R) is the correct explanation of Assertion (A)
  \item (B) Both Assertion (A) and Reason (R) are true but Reason (R) is NOT the correct explanation of Assertion (A)
  \item (C) Assertion (A) is true but Reason (R) is false
  \item (D) Assertion (A) is false but Reason (R) is true
\end{enumerate}
\vspace{0.15cm}

\question[1] Assertion (A): The value of $\sin \theta$ can be $\frac{4}{3}$. Reason (R): The value of $\sin \theta$ is always less than or equal to 1.
\begin{enumerate}[label=(\Alph*)]
  \item (A) Both Assertion (A) and Reason (R) are true and Reason (R) is the correct explanation of Assertion (A)
  \item (B) Both Assertion (A) and Reason (R) are true but Reason (R) is NOT the correct explanation of Assertion (A)
  \item (C) Assertion (A) is true but Reason (R) is false
  \item (D) Assertion (A) is false but Reason (R) is true
\end{enumerate}
\vspace{0.15cm}

\question[1] Assertion (A): If $\tan \theta = 1$, then $\theta = 45^\circ$. Reason (R): $\tan \theta = 1$ only for $\theta = 45^\circ$ in the range $0^\circ \leq \theta < 90^\circ$.
\begin{enumerate}[label=(\Alph*)]
  \item (A) Both Assertion (A) and Reason (R) are true and Reason (R) is the correct explanation of Assertion (A)
  \item (B) Both Assertion (A) and Reason (R) are true but Reason (R) is NOT the correct explanation of Assertion (A)
  \item (C) Assertion (A) is true but Reason (R) is false
  \item (D) Assertion (A) is false but Reason (R) is true
\end{enumerate}
\vspace{0.15cm}

\question[1] Assertion (A): The value of $\cos^2 30^\circ + \sin^2 30^\circ$ is 1. Reason (R): $\sin^2 \theta + \cos^2 \theta = 1$ for all $\theta$.
\begin{enumerate}[label=(\Alph*)]
  \item (A) Both Assertion (A) and Reason (R) are true and Reason (R) is the correct explanation of Assertion (A)
  \item (B) Both Assertion (A) and Reason (R) are true but Reason (R) is NOT the correct explanation of Assertion (A)
  \item (C) Assertion (A) is true but Reason (R) is false
  \item (D) Assertion (A) is false but Reason (R) is true
\end{enumerate}
\vspace{0.15cm}

\question[1] Assertion (A): As $\theta$ increases from $0^\circ$ to $90^\circ$, the value of $\sin \theta$ increases. Reason (R): $\sin 0^\circ = 0$ and $\sin 90^\circ = 1$.
\begin{enumerate}[label=(\Alph*)]
  \item (A) Both Assertion (A) and Reason (R) are true and Reason (R) is the correct explanation of Assertion (A)
  \item (B) Both Assertion (A) and Reason (R) are true but Reason (R) is NOT the correct explanation of Assertion (A)
  \item (C) Assertion (A) is true but Reason (R) is false
  \item (D) Assertion (A) is false but Reason (R) is true
\end{enumerate}
\vspace{0.15cm}

\question[1] Assertion (A): $\sec^2 \theta - 1 = \tan^2 \theta$. Reason (R): $\sec^2 \theta + \tan^2 \theta = 1$ is a standard identity.
\begin{enumerate}[label=(\Alph*)]
  \item (A) Both Assertion (A) and Reason (R) are true and Reason (R) is the correct explanation of Assertion (A)
  \item (B) Both Assertion (A) and Reason (R) are true but Reason (R) is NOT the correct explanation of Assertion (A)
  \item (C) Assertion (A) is true but Reason (R) is false
  \item (D) Assertion (A) is false but Reason (R) is true
\end{enumerate}
\vspace{0.15cm}

\end{questions}
\bigskip
\noindent\textbf{\large Section B: Very Short Answer (2 marks each)}

\begin{questions}
\question[2] If $\tan A = \frac{3}{4}$, find the value of $\sin A \cos A$.
\vspace{0.15cm}

\question[2] Evaluate: $\frac{\tan 65^\circ}{\cot 25^\circ} + \sin^2 30^\circ$.
\vspace{0.15cm}

\question[2] If $4 \tan \theta = 3$, evaluate $\frac{4 \sin \theta - \cos \theta}{4 \sin \theta + \cos \theta}$.
\vspace{0.15cm}

\question[2] If $\cos(A + B) = 0$ and $\sin(A - B) = \frac{1}{2}$, find the values of $A$ and $B$, where $0^\circ < A + B \leq 90^\circ$ and $A > B$.
\vspace{0.15cm}

\question[2] Given $\sec \theta = \frac{13}{5}$, calculate $\cot \theta$.
\vspace{0.15cm}

\question[2] If $\tan A = \frac{4}{3}$, find the value of $\sin A + \cos A$.
\vspace{0.15cm}

\question[2] Evaluate $\frac{2 \tan^2 45^\circ + \cos^2 30^\circ - \sin^2 60^\circ}{1}$.
\vspace{0.15cm}

\end{questions}
\bigskip
\noindent\textbf{\large Section C: Short Answer (3 marks each)}

\begin{questions}
\question[3] If $\sin(A + B) = 1$ and $\cos(A - B) = \frac{\sqrt{3}}{2}$, where $0^\circ < A + B \leq 90^\circ$ and $A > B$, find the values of $A$ and $B$.
\vspace{0.15cm}

\question[3] Evaluate the following expression: $\frac{5 \cos^2 60^\circ + 4 \sec^2 30^\circ - \tan^2 45^\circ}{\sin^2 30^\circ + \cos^2 30^\circ}$
\vspace{0.15cm}

\question[3] If $15 \cot A = 8$, find the values of $\sin A$ and $\sec A$.
\vspace{0.15cm}

\end{questions}
\bigskip
\noindent\textbf{\large Section D: Long Answer (4-5 marks each)}

\begin{questions}
\question[5] Prove that $\frac{\tan \theta}{1 - \cot \theta} + \frac{\cot \theta}{1 - \tan \theta} = 1 + \sec \theta \operatorname{cosec} \theta$. Hence, or otherwise, evaluate the expression if $\tan \theta = 2$.
\vspace{0.15cm}

\question[5] A vertical tower stands on a horizontal plane and is surmounted by a vertical flagstaff of height $6$ m. At a point on the plane, the angle of elevation of the bottom of the flagstaff is $30^\circ$ and that of the top of the flagstaff is $60^\circ$. Find the height of the tower.
\vspace{0.15cm}

\end{questions}
\bigskip

\end{document}
